% =====================================================================
% MCBOMs Methodology
% Mathematical Formulation and Validation
%
% Standalone professional document version of the MCBOMs methodology.
% Formal version of this content (with chapter context) is at
% chapter2.tex / chapter2.pdf — that file is the Task 4 Report deliverable.
% This file is the standalone version for distribution and the project
% website.
%
% Typesetting conventions:
%   - All equations use proper math mode
%   - Variables are italicized; subscripts/superscripts use braces
%   - Numbered equations use \begin{equation} ... \end{equation}
%   - Cross-references use \ref{} for sections, eq., tables
%   - Citations use natural language (Author, Year)
%
% Note on numbering: equations and sections use 2.x notation throughout.
% This is preserved from the source to keep cross-references in the rest
% of the project (AMPL/GAMS/LP files, Python source, USER_MANUAL, GitHub
% Pages site) consistent. The 2.x prefix should be read as a stable
% identifier, not a chapter reference.
%
% Date: April 26, 2026
% Authors: A. Sameen and team, Texas A&M Transportation Institute
% =====================================================================

\documentclass[12pt,letterpaper]{article}

\usepackage[margin=1in]{geometry}
\usepackage{amsmath,amssymb}
\usepackage{booktabs}
\usepackage{array}
\usepackage{tabularx}
\usepackage{enumitem}
\usepackage{hyperref}
\usepackage{xcolor}
\usepackage{siunitx}
\usepackage{caption}
\usepackage{float}

\sisetup{
  group-separator = {,},
  group-minimum-digits = 4,
  output-decimal-marker = {.},
}

\hypersetup{
  colorlinks=false,
  pdfborder={0 0 0},
}

% Custom commands for frequently-used notation
\newcommand{\Btot}{B^{\mathrm{total}}}
\newcommand{\Bsafety}{B^{\mathrm{safety}}}
\newcommand{\Bops}{B^{\mathrm{operations}}}
\newcommand{\Bccm}{B^{\mathrm{ccm}}}
\newcommand{\Enobuild}{E^{\mathrm{no\text{-}build}}}
\newcommand{\Ebuild}{E^{\mathrm{build}}}
\newcommand{\UC}{\mathit{UC}}
\newcommand{\DF}{\mathit{DF}}
\newcommand{\PWF}{\mathit{PWF}}
\newcommand{\CMF}{\mathit{CMF}}
\newcommand{\AADT}{\mathit{AADT}}
\newcommand{\OCC}{\mathit{OCC}}
\newcommand{\VOT}{\mathit{VOT}}
\newcommand{\VOC}{\mathit{VOC}}
\newcommand{\VMT}{\mathit{VMT}}

% Section numbering: preserved from the source document for cross-reference
% consistency with the rest of the MCBOMs project (solver-language model files,
% Python comments, etc.). Numbers 2.x are stable identifiers.
\renewcommand{\thesection}{2.\arabic{section}}
\renewcommand{\thesubsection}{2.\arabic{section}.\arabic{subsection}}
\renewcommand{\thesubsubsection}{2.\arabic{section}.\arabic{subsection}.\arabic{subsubsection}}
\renewcommand{\theequation}{2.\arabic{equation}}
\renewcommand{\thetable}{2.\arabic{table}}
\renewcommand{\thefigure}{2.\arabic{figure}}

\title{MCBOMs Methodology \\ \large Mathematical Formulation and Validation}
\author{Texas A\&M Transportation Institute \\ FHWA Contract HRSO20240009PR}
\date{April 26, 2026}

\begin{document}

\maketitle

\tableofcontents
\newpage

% =====================================================================
% BEGIN MCBOMS METHODOLOGY CONTENT
% =====================================================================

\section*{Introduction}
\addcontentsline{toc}{section}{Introduction}

This document presents the mathematical formulation and validation of the
Multidisciplinary Cost-Benefit Optimization Models (MCBOMs) framework, a
Mixed-Integer Linear Programming (MILP) approach to transportation
infrastructure investment decisions. Given a portfolio of candidate projects
(sites, intersections, or corridors), each with multiple possible improvement
alternatives, MCBOMs selects the combination of alternatives that maximizes
total societal benefit subject to a budget constraint. The benefits include
safety improvements (reductions in expected crash costs), operational
improvements (reductions in travel time and vehicle operating costs), and
corridor condition improvements (energy, environmental, accessibility, and
resilience measures).

The framework is developed under FHWA Contract HRSO20240009PR by the Texas
A\&M Transportation Institute. It builds directly on Harwood, Rabbani, and
Richard~(2003), who developed the Resurfacing Safety Resource Allocation
Program (RSRAP) for allocating resources to safety improvements within
Resurfacing, Restoration, and Rehabilitation (RRR) projects on nonfreeway
facilities, and on Banihashemi~(2007), who extended that framework with
intersection-level optimization using the IHSDM Crash Prediction Module.
Default unit costs and economic parameters come from USDOT Benefit-Cost
Analysis Guidance for Discretionary Grant Programs (May 2025).

This document is organized as follows. Section~2.1 establishes the
foundational framework, including the project--alternative structure, the
benefit definition, and the discounting convention. Section~2.2 presents
the MILP formulation at the project and network levels. Sections~2.3
through~2.5 derive the safety, operational, and corridor condition benefit
components. Section~2.6 documents the solution methodology, and Section~2.7
presents validation against the Harwood~(2003) and Banihashemi~(2007)
benchmarks together with cross-validation against an internal benefit-cost
analysis spreadsheet.

% ---------------------------------------------------------------------
\setcounter{section}{0}
\section{Foundational Framework}
% ---------------------------------------------------------------------

The MCBOMs framework addresses the multi-objective resource allocation problem
faced by transportation agencies: given a portfolio of candidate projects, each
with multiple possible improvement alternatives, determine the combination of
alternatives that maximizes total societal benefit subject to a budget
constraint. Three foundational design choices govern the formulation: the
project--alternative structure, the benefit definition and quantification
structure, and the discounting and time-horizon convention.

% ---------------------------------------------------------------------
\subsection{Project--Alternative Structure}
% ---------------------------------------------------------------------

A \emph{project} represents a candidate location for investment. Examples
include a corridor segment, an intersection, a bridge, or any other discrete
asset for which an improvement decision can be made. Each project is denoted by
the index $i \in \{1, 2, \ldots, I\}$, where $I$ is the total number of
projects in the analysis.

For each project $i$, the analyst defines a set of \emph{improvement
alternatives}, denoted $A_i = \{0, 1, 2, \ldots, J_i\}$. The alternative
$j = 0$ is always the do-nothing baseline (no improvement implemented at site
$i$). Alternatives $j \geq 1$ each represent a distinct combination of
treatments. Treatments may include, for example, lane widening, shoulder paving,
intersection signalization, or pavement resurfacing. A single alternative may
bundle multiple treatments, in which case the alternative's combined effect is
computed using the appropriate combination rule for each benefit category (see
Section~\ref{sec:safety} for safety, Section~\ref{sec:ops} for operations).

Each project is characterized by a record:
\begin{equation}
P_i = \{ \mathrm{Location}_i,\ \mathrm{Cost}_{ij},\ \mathrm{SafetyData}_{ij},\
         \mathrm{OperationData}_{ij},\ \mathrm{CCMData}_{ij},\
         \mathrm{TreatmentType}_{ij} \}
\end{equation}
\noindent for each alternative $j \in A_i$, capturing the data required to
evaluate that alternative.

The optimization decision variable is binary:
\begin{equation}
x_{ij} \in \{0, 1\}, \qquad
x_{ij} = \begin{cases}
1 & \text{if alternative } j \text{ is selected for project } i \\
0 & \text{otherwise}
\end{cases}
\end{equation}

The mutual-exclusivity constraint
$\sum_{j \in A_i} x_{ij} \leq 1$ ensures that at most one alternative (including
do-nothing) is selected for each project.

% ---------------------------------------------------------------------
\subsection{Benefit Definition and Quantification Structure}
% ---------------------------------------------------------------------

The total present-value societal benefit of selecting alternative $j$ for
project $i$, evaluated over the analysis horizon $T$, is the discounted sum
of three additive annual benefit components:
\begin{equation}
\Btot_{ij}\ =\ \sum_{t=1}^{T}\bigl[\, B_{ijt}^{\mathrm{safety}} + B_{ijt}^{\mathrm{operations}} + B_{ijt}^{\mathrm{corridor}} \,\bigr] \cdot \DF_t
\label{eq:btot}
\end{equation}

\noindent where:
\begin{itemize}[leftmargin=2em, itemsep=2pt]
  \item $B_{ijt}^{\mathrm{safety}}$ is the annual monetized safety benefit
    in year $t$ (crash cost reduction); see Section~\ref{sec:safety}.
  \item $B_{ijt}^{\mathrm{operations}}$ is the annual monetized operational
    benefit in year $t$ (travel time, vehicle operating costs, reliability);
    see Section~\ref{sec:ops}.
  \item $B_{ijt}^{\mathrm{corridor}}$ is the annual monetized corridor
    condition benefit in year $t$ (energy, environmental, resilience,
    accessibility); see Section~\ref{sec:ccm}.
  \item $\DF_t = 1/(1+r)^t$ is the discount factor in year $t$, with $r$ the
    real discount rate (Section~\ref{sec:discounting}).
\end{itemize}

For notational compactness in the optimization formulation
(Section~\ref{sec:milp}), the present-value-aggregated forms are denoted
$\Bsafety_{ij}$, $\Bops_{ij}$, and $\Bccm_{ij}$, with the relationship:
\begin{equation*}
\Btot_{ij}\ =\ \Bsafety_{ij} + \Bops_{ij} + \Bccm_{ij}
\end{equation*}

The additive decomposition assumes that the three benefit categories are
independent of one another. Where overlap exists (for example, when a composite
metric such as the Mobility Energy Productivity (MEP) index embeds both travel
time and energy components), explicit double-counting prevention rules are
applied; see Section~\ref{sec:ccm} for the relevant adjustments.

Each project--alternative pair has an associated cost:
\begin{equation}
C_{ij} = C_{ij}^{\mathrm{construction}} + C_{ij}^{\mathrm{indirect}}
\end{equation}
\noindent where construction cost includes resurfacing and geometric
improvements as applicable, and indirect cost captures traffic management
during construction, environmental mitigation, and utility relocation. For the
Task~4 prototype, indirect cost is folded into construction cost where data
permits; otherwise it is set to zero.

Section~\ref{sec:alt-composition} extends this cost definition to the
segment level, showing how $C_{ij}$ is composed from per-treatment
per-segment cost contributions.

% ---------------------------------------------------------------------
\subsection{Alternative Composition and Segment-Level Treatment Specification}
\label{sec:alt-composition}
% ---------------------------------------------------------------------

A central modeling concept in the MCBOMs framework, requested by FHWA
reviewers and consistent with Banihashemi~(2007), is that an alternative
$j$ for project $i$ is not necessarily a single uniform treatment applied
across the entire project. Rather, an alternative is a \emph{composition}
of segment-specific treatments: different homogeneous segments within the
same project may receive different treatment combinations under the same
alternative.

\paragraph{Homogeneous segmentation.} Each project $i$ is partitioned into
$K_i$ homogeneous segments, where a segment $k \in \{1, 2, \ldots, K_i\}$
is a length of facility over which the relevant geometric, traffic, and
crash-history attributes are approximately constant. New segments begin
where the AADT, lane configuration, shoulder treatment, horizontal/vertical
alignment, or roadside condition changes materially. This segmentation rule
matches the Highway Safety Manual (AASHTO, 2010, Part~C) homogeneous-segment
definition and the construction Banihashemi (2007, page~112) uses for his
Crash Prediction Module.

\paragraph{Treatment sets.} For each alternative $j \in A_i$ and each
segment $k \in \{1, \ldots, K_i\}$, the analyst defines a treatment set:
\begin{equation*}
T_{jk}\ =\ \{\,\text{treatments applied to segment } k \text{ under alternative } j\,\}
\end{equation*}

\noindent For the do-nothing alternative ($j = 0$), $T_{0k} = \emptyset$ for
every segment $k$. For non-trivial alternatives, $T_{jk}$ may differ across
segments within the same alternative; this is the segment-level composition
property.

\paragraph{Alternative as a composition.} An alternative $j$ for project $i$
is fully specified by the collection of its segment-level treatment sets:
\begin{equation*}
\text{Alternative } j \text{ for project } i\ =\ \{\, T_{j1},\ T_{j2},\ \ldots,\ T_{jK_i} \,\}
\end{equation*}

\paragraph{Cost decomposition.} The total cost of alternative $j$ for
project $i$ is the sum of treatment-level costs across all segments:
\begin{equation}
C_{ij}\ =\ \sum_{k=1}^{K_i} \sum_{n \in T_{jk}} \mathit{Cost}_{ikn}
\label{eq:cost-decomp}
\end{equation}

\noindent where $\mathit{Cost}_{ikn}$ is the implementation cost of
treatment $n$ on segment $k$ of project $i$.

\paragraph{Benefit decomposition.} Total present-value benefit
$\Btot_{ij}$, defined in Equation~\ref{eq:btot}, similarly aggregates
segment-level contributions. Each annual benefit component
($B_{ijt}^{\mathrm{safety}}$, $B_{ijt}^{\mathrm{operations}}$,
$B_{ijt}^{\mathrm{corridor}}$) is computed as a sum over segments
$k = 1, \ldots, K_i$, with the segment-level value depending on the
treatment set $T_{jk}$ applied to that segment. The detailed
segment-level expansions are given in Section~\ref{sec:safety}
(Equation~\ref{eq:safety-benefit}) for safety, Section~\ref{sec:ops}
(Equation~\ref{eq:ops-benefit}) for operations, and Section~\ref{sec:ccm}
(Equation~\ref{eq:ccm-benefit}) for corridor condition.

\paragraph{Optimization decision variable remains alternative-level.}
Although costs and benefits are computed at the treatment-segment level,
the optimization decision variable $x_{ij}$ remains binary at the
alternative level. The optimizer selects an entire alternative (with its
full collection of segment-specific treatments) or rejects it; partial
alternatives are not selectable. This matches Harwood~(2003)'s site-level
decision-variable structure and is consistent with the Banihashemi~(2007)
formulation interpreted at the alternative level.

\paragraph{Illustrative example.} Consider a 6-mile rural two-lane highway
project ($i = 1$) partitioned into three 2-mile homogeneous segments
($K_1 = 3$). A single non-trivial alternative ($j = 1$) is defined as the
following composition:
\begin{itemize}[leftmargin=2em, itemsep=2pt]
  \item Segment $k = 1$ (miles 0--2): left-turn lane addition at the
    intersection located within this segment. $T_{11} = \{\,\text{left-turn lane}\,\}$.
  \item Segment $k = 2$ (miles 2--4): centerline rumble strips.
    $T_{12} = \{\,\text{rumble strips}\,\}$.
  \item Segment $k = 3$ (miles 4--6): traffic signal installation at the
    intersection located within this segment.
    $T_{13} = \{\,\text{traffic signal}\,\}$.
\end{itemize}

\noindent The total cost of this alternative is computed by
Equation~\ref{eq:cost-decomp}:
\begin{align*}
C_{11}\ &=\ \mathit{Cost}_{1,1,\text{LTL}} + \mathit{Cost}_{1,2,\text{rumble}} + \mathit{Cost}_{1,3,\text{signal}} \\
       &=\ \$300{,}000 + \$50{,}000 \cdot 2\ \text{mi} + \$250{,}000 \\
       &=\ \$650{,}000
\end{align*}

\noindent and total present-value benefit $\Btot_{11}$ is the sum of
segment-level safety, operational, and corridor benefits computed under
each segment's treatment set, discounted via Equation~\ref{eq:btot}.

\paragraph{Combined-treatment effects within a segment.} When a single
segment $k$ receives multiple treatments simultaneously
($\lvert T_{jk} \rvert > 1$), the combined safety effect uses the
multiplicative Crash Modification Factor rule
(Equation~\ref{eq:cmf-combine}). For example, a segment receiving both
shoulder widening (CMF = $0.93$) and rumble strips (CMF = $0.86$) has a
combined CMF of $0.93 \times 0.86 = 0.80$. This rule is the HSM-recommended
default; ongoing CMF Clearinghouse research is refining combined-treatment
estimates.

\paragraph{Implementation note.} The Task~4 prototype dataset encodes
alternative composition through the alternative identifier itself
(e.g., ``Pavement Marking, Cable Barrier'' as a single alternative label).
A treatment-level cost catalog is required to evaluate
Equation~\ref{eq:cost-decomp} term-by-term; in the prototype's BCA
spreadsheet, this catalog is partially captured through per-mile unit cost
fields by treatment type. Full treatment-level decomposition (in which
each $\mathit{Cost}_{ikn}$ is computed from a separate treatment-cost table
joined to a segment-treatment assignment table) is documented as a planned
enhancement; the prototype currently uses alternative-level aggregate
costs as an equivalent simplification.

% ---------------------------------------------------------------------
\subsection{Discounting and Analysis Horizon}
\label{sec:discounting}
% ---------------------------------------------------------------------

All benefit and cost streams are discounted to present value using the
USDOT-recommended real discount rate of 7~percent per year (USDOT, 2025,
Section~4.3). The discount factor for year $t$ is:
\begin{equation}
\DF_t = \frac{1}{(1 + r)^t}, \qquad r = 0.07
\label{eq:df}
\end{equation}

For benefit streams that are constant across years (the standard assumption for
the Task~4 prototype), the present-value calculation reduces to multiplication
by the uniform-series present worth factor:
\begin{equation}
\PWF(r, T) = \sum_{t=1}^{T} \DF_t = \frac{(1+r)^T - 1}{r \cdot (1+r)^T}
\label{eq:pwf}
\end{equation}

\noindent At the framework defaults of $r = 0.07$ and $T = 20$ years, this
evaluates to $\PWF(0.07, 20) = 10.594$.

\paragraph{Analysis horizon.} The Task~4 prototype uses a uniform 20-year
analysis horizon for all alternatives, consistent with USDOT BCA Guidance
(May~2025, Section~4.4) for capacity and operating improvements and with
Banihashemi (2007), who applied a uniform 20-year project lifetime to all
improvement alternatives in his case study.

\paragraph{Treatment-specific service life.} Table~\ref{tab:service-life}
documents typical service lives by treatment type for reference. These values
are not currently applied to the prototype's benefit calculations; all benefits
are evaluated over the uniform 20-year horizon. This simplification applies
uniformly across alternatives, which preserves the relative ranking that drives
the optimization. Per-treatment service life and residual-value adjustments
(in the sense of USDOT BCA Guidance Section~6.3, which applies residual value
to construction costs rather than to benefit streams) are documented as
planned enhancements.

\begin{table}[H]
\centering
\caption{Reference service lives by treatment type. These values are
documented for reference and for future model enhancements; they are not
currently applied to the Task~4 prototype's benefit calculations, which use a
uniform 20-year analysis horizon.}
\label{tab:service-life}
\begin{tabular}{lc}
\toprule
\textbf{Treatment Type} & \textbf{Reference Service Life (years)} \\
\midrule
Pavement resurfacing                  & 10 \\
Lane widening                         & 20 \\
Shoulder widening / paving            & 20 \\
Horizontal curve improvement          & 20 \\
Vertical alignment improvement        & 20 \\
Intersection signalization            & 15 \\
Roundabout                            & 25 \\
Turn lane addition                    & 20 \\
Climbing/passing lane                 & 20 \\
Sight distance improvement            & 20 \\
\bottomrule
\end{tabular}
\end{table}

% ---------------------------------------------------------------------
\section{Mixed-Integer Linear Programming Formulation}
\label{sec:milp}
% ---------------------------------------------------------------------

The MCBOMs optimization problem is formulated as a Mixed-Integer Linear
Program. This section presents the project-level formulation, in which only a
single budget constraint applies, followed by the network-level formulation,
which adds facility-type and regional sub-budget constraints.

% ---------------------------------------------------------------------
\subsection{Project-Level Formulation}
\label{sec:milp-project}
% ---------------------------------------------------------------------

The project-level MILP maximizes total net benefit subject to a single
program-wide budget. The formulation consists of an objective function and
three constraint families.

\paragraph{Objective function.}
\begin{equation}
\max\ Z\ =\ \sum_{i=1}^{I} \sum_{j \in A_i} \bigl( \Btot_{ij} - C_{ij}^{\mathrm{disc}} \bigr) \cdot x_{ij}
\label{eq:obj-project}
\end{equation}

The objective expresses total net societal benefit as the sum, over all
projects and selected alternatives, of benefit minus the
\emph{discretionary} cost component $C_{ij}^{\mathrm{disc}}$. The
parentheses around $(\Btot_{ij} - C_{ij}^{\mathrm{disc}})$ are essential:
the multiplication by $x_{ij}$ applies to the entire net-benefit term,
not to the cost alone.

\paragraph{Discretionary versus committed cost.} For projects where
project cost decomposes into a committed component $C_{ij}^{\mathrm{comm}}$
(expenditure that will occur regardless of the optimization outcome,
typically because policy or asset condition mandates it) and a
discretionary component $C_{ij}^{\mathrm{disc}}$ (expenditure subject to
optimization choice), the objective subtracts only the discretionary
portion. The total cost remains $C_{ij} = C_{ij}^{\mathrm{comm}} +
C_{ij}^{\mathrm{disc}}$ and is the quantity used in the budget
constraint (Equation~\ref{eq:budget-project}).

This distinction is consistent with Harwood et al.~(2003), who model
RRR projects in which pavement resurfacing is committed (every site is
resurfaced under the program) and only the safety improvement cost is
optimized. In that case:
\begin{align*}
C_{ij}^{\mathrm{comm}} &= \text{resurfacing cost (committed by the RRR program)} \\
C_{ij}^{\mathrm{disc}} &= \text{safety improvement cost (the optimization variable)}
\end{align*}

For projects where no committed component exists ---
$C_{ij}^{\mathrm{comm}} = 0$ and $C_{ij}^{\mathrm{disc}} = C_{ij}$ ---
Equation~\ref{eq:obj-project} reduces to the simple form
$\max \sum (\Btot_{ij} - C_{ij}) x_{ij}$, which is the case for the
Banihashemi~(2007) validation and for general projects without a fixed
infrastructure baseline.

The MCBOMs implementation supports both conventions: the optimizer
accepts a per-alternative objective coefficient, which the analyst sets
to $(\Btot_{ij} - C_{ij}^{\mathrm{disc}})$ when modeling a Harwood-style
problem and to $(\Btot_{ij} - C_{ij})$ for a generic problem.

\paragraph{Budget constraint.}
\begin{equation}
\sum_{i=1}^{I} \sum_{j \in A_i} C_{ij} \cdot x_{ij}\ \leq\ B^{\mathrm{total}}
\label{eq:budget-project}
\end{equation}

\noindent where $B^{\mathrm{total}}$ is the total program budget.

\paragraph{Mutual exclusivity.}
\begin{equation}
\sum_{j \in A_i} x_{ij}\ \leq\ 1, \qquad \forall i \in \{1, 2, \ldots, I\}
\label{eq:mutex}
\end{equation}

This constraint allows at most one alternative (including do-nothing) to be
selected for each project. When the inequality is strict ($\sum_j x_{ij} = 0$),
the optimizer has implicitly chosen do-nothing for project $i$.

\paragraph{Binary decision variables.}
\begin{equation}
x_{ij} \in \{0, 1\}, \qquad \forall i,\ \forall j \in A_i
\label{eq:binary-project}
\end{equation}

\paragraph{Optional constraints.} Beyond the four core constraints above,
the project-level formulation supports three optional constraint families
that reflect common real-world policy and engineering requirements. These
constraints are activated only when the corresponding data is provided.

\textbf{(a) Dependency constraints.} If alternative $j$ of project $i$
requires the prior completion of alternative $j'$ of project $i'$ (for
example, a corridor signalization plan that depends on an upstream
intersection upgrade), the dependency is enforced by:
\begin{equation}
x_{ij}\ \leq\ x_{i'j'}
\label{eq:dependency}
\end{equation}

\textbf{(b) Cross-project exclusivity.} If alternatives at different
projects conflict (for example, when two adjacent projects compete for
the same right-of-way or land use), the conflict is enforced by:
\begin{equation}
x_{ij} + x_{i'j'}\ \leq\ 1
\label{eq:cross-exclusivity}
\end{equation}

\textbf{(c) Minimum benefit-cost ratio.} To require selected alternatives
to exceed a minimum benefit-cost ratio:
\begin{equation}
\sum_{j \in A_i} \Btot_{ij} \cdot x_{ij}\ \geq\ \theta \cdot \sum_{j \in A_i} C_{ij} \cdot x_{ij}, \qquad \forall i
\label{eq:min-bcr}
\end{equation}

\noindent where $\theta$ is the minimum acceptable benefit-cost ratio
(for example, $\theta = 1.0$ enforces positive net benefit at every
selected project; $\theta = 2.0$ enforces a $2{:}1$ benefit-cost ratio).

\paragraph{Implementation status of optional constraints.} The MCBOMs
prototype optimizer currently implements the four core constraints
(Equations~\ref{eq:obj-project}--\ref{eq:binary-project}); the three
optional constraints are documented in this work as supported by the
mathematical framework and are added to the optimizer as agency-data
becomes available to populate them. The Harwood~(2003) and
Banihashemi~(2007) validation case studies do not exercise dependency or
cross-project exclusivity constraints; minimum-BCR enforcement is
optional and disabled by default.

\paragraph{Facility-type scope.} The project-level optimization framework
applies to the following facility types, with facility-specific Safety
Performance Functions (SPFs), Crash Modification Factors (CMFs), and
operational performance parameters as defined in
Section~\ref{sec:safety} (safety) and Section~\ref{sec:ops} (operations):

\begin{itemize}[leftmargin=2em, itemsep=2pt]
  \item \textbf{Rural highways:} two-lane highways and intersections
    (including roundabouts); four-lane divided and undivided highways
    and intersections; freeways, ramps, and ramp terminals.
  \item \textbf{Urban / suburban facilities:} two-lane highways and
    intersections; multilane highways and intersections; one-way
    arterials and intersections; freeways, ramps, and ramp terminals.
\end{itemize}

% ---------------------------------------------------------------------
\subsection{Network-Level Formulation}
\label{sec:milp-network}
% ---------------------------------------------------------------------

The network-level formulation extends the project-level MILP by adding
facility-type sub-budget constraints, regional cap constraints, and regional
minimum-investment constraints. These additional constraints reflect common
real-world resource allocation requirements, such as ensuring that no single
facility class consumes a disproportionate share of the budget, or that each
geographic region receives a minimum allocation of investment.

\paragraph{Sets for network-level constraints.}
Let $\mathcal{F}_m$ denote the set of project indices belonging to facility
type $m \in \{m_1, m_2, \ldots, m_M\}$ (e.g., rural two-lane, rural multilane,
urban arterial). Let $\mathcal{R}_r$ denote the set of project indices located
in region $r \in \{r_1, r_2, \ldots, r_R\}$ (e.g., District~1, District~2).

\paragraph{Objective function.} Identical to the project-level objective
(Equation~\ref{eq:obj-project}); the summation index $i$ now ranges over all
projects across all facility types and regions.

\paragraph{Total budget constraint.} Identical to
Equation~\ref{eq:budget-project}.

\paragraph{Mutual exclusivity.} Identical to Equation~\ref{eq:mutex}.

\paragraph{Facility-type sub-budget constraints.}
\begin{equation}
\sum_{i \in \mathcal{F}_m} \sum_{j \in A_i} C_{ij} \cdot x_{ij}\ \leq\ B_m,
\qquad \forall m
\label{eq:budget-facility}
\end{equation}

\noindent where $B_m$ is the maximum allocation to facility type $m$.

\paragraph{Regional sub-budget constraints (cap).}
\begin{equation}
\sum_{i \in \mathcal{R}_r} \sum_{j \in A_i} C_{ij} \cdot x_{ij}\ \leq\ B_r,
\qquad \forall r
\label{eq:budget-region-cap}
\end{equation}

\paragraph{Regional minimum-investment constraints (floor).}
\begin{equation}
\sum_{i \in \mathcal{R}_r} \sum_{j \in A_i} C_{ij} \cdot x_{ij}\ \geq\ \beta_r \cdot B^{\mathrm{total}},
\qquad \forall r
\label{eq:budget-region-floor}
\end{equation}

\noindent where $\beta_r \in [0, 1]$ is the minimum-investment fraction for
region $r$. Setting $\beta_r = 0$ disables the floor constraint for that
region.

\paragraph{Binary decision variables.} Identical to
Equation~\ref{eq:binary-project}.

% ---------------------------------------------------------------------
\subsection{Worked Example: Project-Level Optimization for Three Rural Highway Segments}
\label{sec:milp-worked}
% ---------------------------------------------------------------------

To illustrate the project-level MILP, consider three rural two-lane
highway projects with the following alternative inventory and a program
budget of $\$2{,}000{,}000$. Costs and total present-value benefits are
expressed in thousands of dollars.

\textbf{Project 1 (5-mile rural segment):}
\begin{itemize}[leftmargin=2em, itemsep=2pt]
  \item $j = 0$ (do nothing): $C_{10} = 0$, $\Btot_{10} = 0$
  \item $j = 1$ (widen lanes to 12~ft): $C_{11} = \$800$, $\Btot_{11} = \$1{,}200$
  \item $j = 2$ (widen shoulders to 6~ft): $C_{12} = \$600$, $\Btot_{12} = \$950$
\end{itemize}

\textbf{Project 2 (3-mile rural segment with curve):}
\begin{itemize}[leftmargin=2em, itemsep=2pt]
  \item $j = 0$: $C_{20} = 0$, $\Btot_{20} = 0$
  \item $j = 1$ (flatten curve): $C_{21} = \$1{,}500$, $\Btot_{21} = \$2{,}400$
\end{itemize}

\textbf{Project 3 (2-mile rural segment):}
\begin{itemize}[leftmargin=2em, itemsep=2pt]
  \item $j = 0$: $C_{30} = 0$, $\Btot_{30} = 0$
  \item $j = 1$ (rumble strips): $C_{31} = \$100$, $\Btot_{31} = \$150$
  \item $j = 2$ (widen lanes + shoulders): $C_{32} = \$1{,}200$, $\Btot_{32} = \$1{,}800$
\end{itemize}

\paragraph{MILP instance.} The objective function and constraints
specialize to:
\begin{equation*}
\begin{aligned}
\text{Maximize:}\quad & Z = 400 x_{11} + 350 x_{12} + 900 x_{21} + 50 x_{31} + 600 x_{32} \\[4pt]
\text{Subject to:}\quad & 800 x_{11} + 600 x_{12} + 1500 x_{21} + 100 x_{31} + 1200 x_{32}\ \leq\ 2000 \\
& x_{10} + x_{11} + x_{12}\ \leq\ 1 \\
& x_{20} + x_{21}\ \leq\ 1 \\
& x_{30} + x_{31} + x_{32}\ \leq\ 1 \\
& x_{ij} \in \{0, 1\}
\end{aligned}
\end{equation*}

\noindent where the objective coefficients are net benefits
$(\Btot_{ij} - C_{ij})$ for non-zero alternatives. This worked example
has no committed cost component
($C_{ij}^{\mathrm{comm}} = 0$,\ $C_{ij}^{\mathrm{disc}} = C_{ij}$), so
the objective uses the simple form.

\paragraph{Optimal solution.} A MILP solver evaluates all feasible
combinations of project--alternative selections and returns the
combination that maximizes total net benefit subject to the budget
constraint. The optimal solution for this instance is:
\begin{itemize}[leftmargin=2em, itemsep=2pt]
  \item $x_{11} = 1$ (Project~1, lane widening) -- net benefit \$400, cost \$800
  \item $x_{20} = 1$ (Project~2, do nothing) -- net benefit \$0, cost \$0
  \item $x_{32} = 1$ (Project~3, lane + shoulder widening) -- net benefit \$600, cost \$1{,}200
\end{itemize}

\noindent \textbf{Total cost:} $\$800 + \$0 + \$1{,}200 = \$2{,}000$
(at budget). \textbf{Total net benefit:} $\$400 + \$600 = \$1{,}000$
(thousand dollars).

This example illustrates two key features of the MILP formulation:
(a)~the trade-off between high-cost / high-benefit alternatives
(Project~2's curve flattening is rejected despite its $\$900$ net benefit
because its $\$1{,}500$ cost would crowd out higher-net-benefit
combinations elsewhere); and (b)~the role of the do-nothing alternative
in providing a feasible default when budget constraints prevent funding
any improvement at a given site.

% ---------------------------------------------------------------------
\section{Safety Benefit Quantification}
\label{sec:safety}
% ---------------------------------------------------------------------

The safety benefit module quantifies the present-value reduction in expected
crash costs resulting from each improvement alternative. The formulation
follows the methodology of the Highway Safety Manual (AASHTO, 2010, Part~C),
extended to the optimization context developed by Harwood et al.~(2003) and
Banihashemi~(2007), and updated to use current USDOT-recommended unit crash
costs.

% ---------------------------------------------------------------------
\subsection{Severity Disaggregation}
% ---------------------------------------------------------------------

The MCBOMs framework adopts the five-bin KABCO crash severity classification
as the standard:
\begin{itemize}[leftmargin=2em, itemsep=2pt]
  \item \textbf{K} --- Killed (fatal injury)
  \item \textbf{A} --- Suspected serious (incapacitating) injury
  \item \textbf{B} --- Suspected minor (non-incapacitating) injury
  \item \textbf{C} --- Possible injury
  \item \textbf{O} --- Property damage only (no injury)
\end{itemize}

The severity index set is denoted $\mathcal{S} = \{K, A, B, C, O\}$. All
benefit calculations sum over this five-element set.

\paragraph{Variants for benchmark reproduction.} Two legacy severity
conventions are used in the validation reproductions described in
Section~\ref{sec:validation}: Harwood (2003) used a two-bin convention
(fatal+injury combined; PDO), and Banihashemi (2007) collapsed five severity
levels into a single average crash cost per facility type. These variants are
documented as conversion bridges and applied only when reproducing the
respective papers' published results.

% ---------------------------------------------------------------------
\subsection{Unit Crash Costs}
% ---------------------------------------------------------------------

The framework default unit crash costs are taken from USDOT BCA Guidance
(May~2025), Appendix~A, Table~A-1 (page~39), in 2023 dollars:
\begin{equation*}
\UC_K = \$13{,}200{,}000, \quad
\UC_A = \$1{,}254{,}700, \quad
\UC_B = \$246{,}900, \quad
\UC_C = \$118{,}000, \quad
\UC_O = \$5{,}300
\end{equation*}

These values include both economic costs (medical, vehicle damage, lost wages)
and quality-adjusted-life-year (QALY) costs. State agencies applying the
MCBOMs framework may substitute state-specific crash costs in place of the
national defaults; FHWA-SA-25-021 (October~2025), Step~5, provides a
per-capita-income adjustment procedure for converting national values to
state-level values.

% ---------------------------------------------------------------------
\subsection{Safety Benefit Equation}
% ---------------------------------------------------------------------

The present-value safety benefit for project $i$, alternative $j$, over the
analysis horizon $T$ is:
\begin{equation}
\Bsafety_{ij}\ =\ \sum_{t=1}^{T}\sum_{k=1}^{K_i}\sum_{s \in \mathcal{S}}
  \bigl[\Enobuild_{ikst} - \Ebuild_{ijkst}\bigr] \cdot \UC_s \cdot \DF_t
\label{eq:safety-benefit}
\end{equation}

\noindent where:
\begin{itemize}[leftmargin=2em, itemsep=2pt]
  \item $K_i$ is the number of homogeneous segments (or intersections)
    composing project $i$, per HSM Part~C segmentation
  \item $\Enobuild_{ikst}$ is the expected number of crashes of severity $s$
    on segment $k$ in year $t$ under the no-build (existing) condition
  \item $\Ebuild_{ijkst}$ is the corresponding expected count under the build
    condition (alternative $j$ implemented)
  \item $\UC_s$ is the unit crash cost for severity $s$
  \item $\DF_t$ is the discount factor from Equation~\ref{eq:df}
\end{itemize}

The build-condition expected crash count is obtained by applying the relevant
Crash Modification Factor (CMF):
\begin{equation}
\Ebuild_{ijkst}\ =\ \Enobuild_{ikst} \cdot \CMF_{ijks}
\label{eq:build-crashes}
\end{equation}

\noindent For alternatives that bundle multiple individual treatments
$n = 1, 2, \ldots, N_{ij}$, the combined CMF is computed using the
HSM-recommended multiplicative rule:
\begin{equation}
\CMF^{\mathrm{combined}}_{ij}\ =\ \prod_{n=1}^{N_{ij}} \CMF_n
\label{eq:cmf-combine}
\end{equation}

This rule assumes that individual treatment effects are independent. The HSM
acknowledges this as a known simplification; ongoing research in the FHWA CMF
Clearinghouse is refining combined CMF estimates, but no consensus alternative
exists at the time of this writing.

% ---------------------------------------------------------------------
\subsection{Constant-Annual-Crashes Simplification}
% ---------------------------------------------------------------------

For the Task~4 prototype, the no-build expected crash count is assumed
constant across years (no traffic growth or treatment-effectiveness decay):
\begin{equation*}
\Enobuild_{ikst} = \Enobuild_{iks}, \qquad \forall t \in \{1, 2, \ldots, T\}
\end{equation*}

Under this assumption, Equation~\ref{eq:safety-benefit} factors as:
\begin{equation}
\Bsafety_{ij}\ =\ \biggl[ \sum_{k=1}^{K_i} \sum_{s \in \mathcal{S}}
  \bigl(\Enobuild_{iks} - \Ebuild_{ijks}\bigr) \cdot \UC_s \biggr]
  \cdot \PWF(r, T)
\label{eq:safety-benefit-pwf}
\end{equation}

This is the form implemented in the MCBOMs Python module
\texttt{mcboms.benefits.safety}. Time-varying expectations (for traffic growth
or effectiveness decay) are supported by an alternative function
(\texttt{safety\_benefit\_time\_varying}) but are not exercised in the Task~4
prototype.

% ---------------------------------------------------------------------
\subsection{Source of $\Enobuild$ for the Prototype}
% ---------------------------------------------------------------------

For the Task~4 prototype, the no-build expected crash count
$\Enobuild_{iks}$ is drawn directly from observed crash counts (typically a
3- to 5-year crash history annualized). This convention matches:
\begin{itemize}[leftmargin=2em, itemsep=2pt]
  \item Banihashemi~(2007), who explicitly states: ``This optimization model
    is for cases in which the EB procedure is not appropriate or there are
    not enough crash data to support this process'' (page~6).
  \item Harwood~(2003), who uses observed crash counts subdivided by default
    severity proportions (page~4).
  \item The parallel data-curation BCA spreadsheet developed in Task~4, which
    uses observed crashes directly.
\end{itemize}

Empirical Bayes adjustment of $\Enobuild$ per Hauer~(1997) is a documented
methodological option for implementations with adequate crash history data.
EB adjustment, when applied, operates as a preprocessing step that produces
an EB-adjusted $\Enobuild$ consumed unchanged by
Equation~\ref{eq:safety-benefit}. The MCBOMs framework neither requires nor
forbids EB adjustment; it is a deployment-time choice.

% ---------------------------------------------------------------------
\subsection{Connection to the Optimization Decision Variable}
% ---------------------------------------------------------------------

The safety benefit $\Bsafety_{ij}$ defined in Equation~\ref{eq:safety-benefit}
is computed during preprocessing for every alternative $j$ of every project
$i$. The optimizer's objective function
(Equation~\ref{eq:obj-project}) multiplies $\Bsafety_{ij}$ by the binary
decision variable $x_{ij}$, ensuring that benefits are credited only when an
alternative is actually selected.

The expected crash count that materializes after the optimization completes
is given by the gating identity:
\begin{equation}
E_{iks}^{\mathrm{final}}\ =\ \biggl(1 - \sum_{j \in A_i} x_{ij}\biggr) \Enobuild_{iks}\ +\ \sum_{j \in A_i} x_{ij} \cdot \Ebuild_{ijks}
\label{eq:gating}
\end{equation}

\noindent with $\sum_{j \in A_i} x_{ij} \leq 1$. This identity is included
for expository completeness; it is not needed in computation, since
$\Bsafety_{ij}$ is precomputed and the optimizer operates on
Equation~\ref{eq:obj-project} directly.

% ---------------------------------------------------------------------
\subsection{Worked Example}
\label{sec:safety-worked-example}
% ---------------------------------------------------------------------

To illustrate the application of Equation~\ref{eq:safety-benefit-pwf},
consider a single rural two-lane segment with the following characteristics:

\begin{itemize}[leftmargin=2em, itemsep=2pt]
  \item Segment length: 5 miles
  \item AADT: 5{,}000 vehicles per day
  \item Observed annual crash frequency $\Enobuild = 6.0$ crashes per year
  \item Severity distribution (from observed crash history):
    1.3\% K, 6.5\% A, 9.5\% B, 11.7\% C, 71.0\% O
  \item Proposed treatment: lane widening from 10~ft to 12~ft, with a single
    CMF of $0.88$ applied uniformly across all severities
  \item Discount rate $r = 0.07$, analysis horizon $T = 20$ years
\end{itemize}

\paragraph{Step 1: Disaggregate observed crashes by severity.}
\[
\Enobuild_K = 0.078,\quad
\Enobuild_A = 0.390,\quad
\Enobuild_B = 0.570,\quad
\Enobuild_C = 0.702,\quad
\Enobuild_O = 4.260
\]
\quad (crashes per year)

\paragraph{Step 2: Apply the CMF to obtain build-condition expected crashes.}
With $\CMF = 0.88$:
\[
\Ebuild_K = 0.0686,\quad
\Ebuild_A = 0.3432,\quad
\Ebuild_B = 0.5016,\quad
\Ebuild_C = 0.6178,\quad
\Ebuild_O = 3.7488
\]

\paragraph{Step 3: Compute annual benefit (sum over severity).}
The reduction is $(1 - \CMF) = 0.12$ of each severity bin. Multiplying by the
USDOT BCA May~2025 unit costs:
\begin{align*}
\text{Annual safety benefit} \ &=\ (0.12)(0.078)(\$13{,}200{,}000) \\
&+ (0.12)(0.390)(\$1{,}254{,}700) \\
&+ (0.12)(0.570)(\$246{,}900) \\
&+ (0.12)(0.702)(\$118{,}000) \\
&+ (0.12)(4.260)(\$5{,}300) \\
&=\ \$123{,}552 + \$58{,}720 + \$16{,}888 + \$9{,}940 + \$2{,}709 \\
&=\ \$211{,}809 \ \text{per year}
\end{align*}

\paragraph{Step 4: Convert to present value.}
\begin{equation*}
\Bsafety = \$211{,}809 \times \PWF(0.07, 20) = \$211{,}809 \times 10.594 = \$2{,}243{,}913
\end{equation*}

This worked example uses the framework default unit crash costs (USDOT BCA
May~2025). The MCBOMs Python implementation reproduces this calculation
exactly; see Section~\ref{sec:validation-worked} for the validation result.

% ---------------------------------------------------------------------
\section{Operational Benefit Quantification}
\label{sec:ops}
% ---------------------------------------------------------------------

The operational benefit module quantifies the present-value reduction in
travel time and vehicle operating costs resulting from each improvement
alternative. The formulation follows USDOT BCA Guidance (May~2025,
Section~5.2 and Appendix~A Tables~A-2 and~A-3) and generalizes the
operational-benefit treatments in Harwood (2003)'s PTOB term and
Banihashemi (2007)'s delay-cost term, both of which are recovered as special
cases of the MCBOMs formulation.

% ---------------------------------------------------------------------
\subsection{Operational Benefit Equation}
% ---------------------------------------------------------------------

The present-value operational benefit for project $i$, alternative $j$, is:
\begin{equation}
\Bops_{ij}\ =\ \sum_{t=1}^{T}\sum_{k=1}^{K_i} \biggl[
  \sum_{v \in \mathcal{V}} \Delta d_{ijkvt} \cdot \AADT_{ijkvt} \cdot 365 \cdot \OCC_v \cdot \VOT_v
  \ +\ \sum_{v \in \mathcal{V}} \Delta\VOC_{ijkv} \cdot \VMT_{ijkvt}
\biggr] \cdot \DF_t
\label{eq:ops-benefit}
\end{equation}

\noindent where:
\begin{itemize}[leftmargin=2em, itemsep=2pt]
  \item $\Delta d_{ijkvt}$ is the change in per-vehicle delay (hours per
    vehicle), for vehicle class $v$ on segment $k$ in year $t$, attributable
    to alternative $j$ at project $i$. Computed externally via Highway
    Capacity Manual procedures, IHSDM, TWOPAS, or microsimulation.
  \item $\AADT_{ijkvt}$ is the annual average daily traffic by vehicle class
    (vehicles per day)
  \item $365$ is the annualization factor (days per year)
  \item $\OCC_v$ is the average vehicle occupancy by vehicle class (persons
    per vehicle), drawn from USDOT BCA Guidance Appendix~A Table~A-3
  \item $\VOT_v$ is the value of travel time by vehicle class (\$ per
    person-hour), drawn from USDOT BCA Guidance Appendix~A Table~A-2
  \item $\Delta\VOC_{ijkv}$ is the change in vehicle operating cost (\$ per
    vehicle-mile)
  \item $\VMT_{ijkvt}$ is the vehicle-miles traveled per year, by segment and
    vehicle class
  \item $\DF_t$ is the discount factor from Equation~\ref{eq:df}
  \item $\mathcal{V}$ is the set of vehicle classes (e.g., light-duty
    passenger, single-unit truck, combination truck)
\end{itemize}

\paragraph{Dimensional consistency.} The travel time term has units
$(\text{hr}/\text{veh}) \cdot (\text{veh}/\text{day}) \cdot (\text{days}) \cdot (\text{persons}/\text{veh}) \cdot (\$/\text{person-hr}) = \$$.
The vehicle operating cost term has units
$(\$/\text{veh-mi}) \cdot (\text{veh-mi}) = \$$.

\paragraph{Standard parameter values.} For the Task~4 prototype, the
following USDOT BCA Guidance May~2025 values apply by default:
\begin{itemize}[leftmargin=2em, itemsep=2pt]
  \item $\VOT$ (all-purposes blend, passenger): \$21.10 per person-hour
    (Table~A-2, page~40)
  \item $\OCC$ (passenger vehicles, all travel): 1.52 persons per vehicle
    (Table~A-3, page~41)
  \item $\VOC$ (light-duty vehicles): \$0.56 per vehicle-mile (Table~A-4,
    page~41)
  \item $\VOC$ (commercial trucks): \$1.27 per vehicle-mile (Table~A-4)
\end{itemize}

% ---------------------------------------------------------------------
\subsection{Generalization of Harwood and Banihashemi Operational Treatments}
% ---------------------------------------------------------------------

Equation~\ref{eq:ops-benefit} generalizes the operational-benefit treatments
in the two benchmark studies; both are recovered as special cases.

\paragraph{Harwood (2003) PTOB as a special case.} Harwood's
Present-value Traffic Operational Benefit (PTOB) captures the travel time
savings from a 1~mph speed increase observed after pavement resurfacing,
valued at \$10.00 per hour (1994 dollars), assumed to persist for 30~months.
Setting $\OCC_v = 1$, $\VOT_v = \$10.00$/hr, $\Delta\VOC = 0$, and limiting
the temporal sum to 30 months ($T = 2.5$ years) at Harwood's 4~percent
discount rate recovers PTOB exactly.

\paragraph{Banihashemi (2007) delay cost as a special case.}
Banihashemi minimizes the sum of crash and delay costs, with delay valued at
\$12.50 per vehicle-hour (page~109), computed as
$\$10/\text{hr} \times 1.25$ persons per vehicle (occupancy folded into a
single rate). Delay times $D$ are computed externally in vehicle-hours per
year using IHSDM/TWOPAS for highway segments and Highway Capacity Manual
procedures (Equations A30-11 and A30-15) for intersections. Setting
$\OCC_v \cdot \VOT_v = \$12.50$/veh-hr (a single combined rate),
$\Delta\VOC = 0$, and treating $\Delta d \cdot \AADT \cdot 365$ as
Banihashemi's annual $D$ recovers his delay cost term exactly.

This generalization allows the MCBOMs framework to be applied across a wider
range of project types than either benchmark addresses individually, while
remaining consistent with both formulations as special cases.

% ---------------------------------------------------------------------
\subsection{Worked Example}
% ---------------------------------------------------------------------

To illustrate the application of Equation~\ref{eq:ops-benefit}, consider a
single highway segment with the following characteristics:

\begin{itemize}[leftmargin=2em, itemsep=2pt]
  \item Segment length: 5 miles
  \item AADT: 1{,}500 vehicles per day (passenger vehicles only)
  \item Per-vehicle delay reduction: $\Delta d = 17$ seconds per vehicle
    $= 17/3600 = 0.00472$ hr per vehicle
  \item Average occupancy $\OCC = 1.52$ persons per vehicle
  \item Value of travel time $\VOT = \$21.10$ per person-hour (USDOT BCA
    Guidance May~2025, all-purposes)
  \item Vehicle operating cost change $\Delta\VOC = 0$ (no change in
    operating cost; the alternative reduces delay only)
  \item Discount rate $r = 0.07$, analysis horizon $T = 20$ years
\end{itemize}

\paragraph{Step 1: Compute annual person-hours of delay reduction.}
\begin{align*}
\text{Annual person-hours saved}
\ &=\ \Delta d \cdot \AADT \cdot 365 \cdot \OCC \\
\ &=\ 0.004722 \cdot 1{,}500 \cdot 365 \cdot 1.52 \\
\ &=\ 3{,}930 \ \text{person-hours per year}
\end{align*}

\paragraph{Step 2: Monetize.}
\begin{equation*}
\text{Annual operational benefit} \ =\ 3{,}930 \times \$21.10 \ =\ \$82{,}919 \ \text{per year}
\end{equation*}

\paragraph{Step 3: Convert to present value.}
\begin{equation*}
\Bops \ =\ \$82{,}919 \times \PWF(0.07, 20) \ =\ \$82{,}919 \times 10.594 \ =\ \$878{,}449
\end{equation*}

This worked example uses the framework default parameters. The vehicle operating
cost contribution is zero in this case; for projects that change vehicle
operating cost (for example, by altering the route's grade or surface type),
the second term in Equation~\ref{eq:ops-benefit} contributes additionally.

% ---------------------------------------------------------------------
\section{Corridor Condition Benefit Quantification}
\label{sec:ccm}
% ---------------------------------------------------------------------

Corridor Condition Measures (CCMs) capture benefit categories that fall
outside the conventional safety and operational scope: energy and emissions
reductions, environmental and health co-benefits, resilience to natural
hazards, accessibility improvements, and changes in property value. The
MCBOMs framework provides a generic computational structure for CCM benefits
and a reference table of monetization rates.

% ---------------------------------------------------------------------
\subsection{CCM Benefit Computation Structure}
% ---------------------------------------------------------------------

The present-value CCM benefit is computed as:
\begin{equation}
\Bccm_{ij}\ =\ \sum_{c \in \mathcal{C}} \sum_{t=1}^{T} \sum_{k=1}^{K_i}
  \Delta Q_{ijkt}^{c} \cdot V_t^{c} \cdot \DF_t
\label{eq:ccm-benefit}
\end{equation}

\noindent where:
\begin{itemize}[leftmargin=2em, itemsep=2pt]
  \item $\mathcal{C}$ is the set of CCM categories applicable to the project
  \item $\Delta Q_{ijkt}^{c}$ is the change in the physical quantity for
    category $c$ on segment $k$ in year $t$ attributable to alternative $j$
    at project $i$ (e.g., metric tons of $\mathrm{NO_x}$, gallons of fuel,
    expected damage from flood events)
  \item $V_t^{c}$ is the monetization rate for category $c$ in year $t$
    (\$ per unit of $Q^c$)
  \item $\DF_t$ is the discount factor from Equation~\ref{eq:df}
\end{itemize}

% ---------------------------------------------------------------------
\subsection{CCM Categories and Reference Sources}
% ---------------------------------------------------------------------

The Task~4 prototype identifies eight CCM categories drawn from the
methodology selected in Task~3. For each, Table~\ref{tab:ccm-sources}
identifies the physical quantity to be measured and the recommended
monetization source.

\begin{table}[H]
\centering
\caption{CCM categories: physical quantities and monetization sources.}
\label{tab:ccm-sources}
\small
\begin{tabular}{p{3.4cm}p{4.0cm}p{6.6cm}}
\toprule
\textbf{Category} & \textbf{Physical Quantity ($\Delta Q^c$)} & \textbf{Monetization Source ($V^c$)} \\
\midrule
Energy / Fuel & Annual fuel consumption change (gallons or BTU) & USDOT BCA Guidance Appendix~A Table~A-4 (vehicle operating cost separation; reflected in $\Bops$ when applicable) \\
\addlinespace
Air Quality (NOx, SOx, PM\textsubscript{2.5}) & Annual emissions change (metric tons) & USDOT BCA Guidance Appendix~A Table~A-6 (page~42) \\
\addlinespace
Resilience to natural hazards & Expected annual damages avoided (\$) or service-loss days reduced & FEMA Standard Economic Values v13 (September~2024) \\
\addlinespace
Accessibility / Mobility Energy Productivity (MEP) & Composite mobility index change (per FDOT/NREL methodology) & Jin et al.~(2023), FDOT BDV29-977-66 \\
\addlinespace
Reduced Emergency Service Delay & Avoided property damage from fire-service delay; ambulance survival rates & FEMA methodology adapted for transportation (USDOT BCA Guidance Section~5.7) \\
\addlinespace
Stormwater Runoff & Reduced runoff volume (gallons) and pollutant loads & Local utility cost data; USDOT BCA Guidance Section~5.7 \\
\addlinespace
Wildlife Crash Reduction & Wildlife-vehicle collisions avoided & USDOT BCA Guidance Section~5.7 (treated within safety methodology where applicable) \\
\addlinespace
Property Value Increase & Net land value change attributable to project (\$, one-time) & Local real estate data; USDOT BCA Guidance Section~5.7 \\
\bottomrule
\end{tabular}
\end{table}

% ---------------------------------------------------------------------
\subsection{Double-Counting Prevention}
% ---------------------------------------------------------------------

The additive decomposition $\Btot_{ij} = \Bsafety_{ij} + \Bops_{ij} +
\Bccm_{ij}$ assumes that the three benefit categories are independent. Three
mechanisms enforce this independence:

\paragraph{(1) Data-level tagging.} Each benefit observation entering the
calculation is tagged with a category indicator
$\delta_{c,ikt} \in \{0, 1\}$ identifying its source category. An observation
contributes to exactly one of $\Bsafety$, $\Bops$, or $\Bccm$.

\paragraph{(2) Equation-level gating.} Each component sums only over its
tagged observations. For example, fuel-consumption changes that are already
reflected in the operational $\VOC$ term of Equation~\ref{eq:ops-benefit}
are not separately counted in the energy CCM category.

\paragraph{(3) Composite-metric adjustment.} When a composite index is used
that embeds multiple component categories (for example, the FDOT/NREL MEP
index, which combines travel time, reliability, and energy with weights
$w_{TT} = 0.40$, $w_{Rel} = 0.25$, $w_{Energy} = 0.35$), the operational and
corridor components are reduced proportionally to avoid duplicating the
embedded contributions:
\begin{align*}
\Bops_{ij} \ &=\ (1 - w_{TT}) \cdot B_{ij}^{TT}
            \ +\ (1 - w_{Rel}) \cdot B_{ij}^{Rel}
            \ +\ B_{ij}^{SD} \\
\Bccm_{ij} \ &=\ (1 - w_{Energy}) \cdot B_{ij}^{Energy}
            \ +\ B_{ij}^{\text{other-CCM}}
\end{align*}

\paragraph{Validation.} The pairwise correlation between benefit components
across all $(i, j)$ pairs is monitored as a diagnostic check; correlations
satisfying $|\rho(B^{c_1}, B^{c_2})| > 0.3$ are flagged for manual review of
potential double-counting.

% ---------------------------------------------------------------------
\section{Solution Methodology}
\label{sec:solution}
% ---------------------------------------------------------------------

This section describes the solution approach for the MCBOMs MILP, the
software implementation, and the relationship between the MILP formulation
and alternative optimization paradigms considered during model selection.

% ---------------------------------------------------------------------
\subsection{Mixed-Integer Linear Programming Formulation Summary}
\label{sec:milp-summary}
% ---------------------------------------------------------------------

The complete MCBOMs MILP is summarized below. The formulation has four
core constraint families that are always active and three optional
constraint families that are activated only when the corresponding agency
data is provided.
\begin{equation*}
\begin{aligned}
\text{Maximize:}\quad & Z = \sum_{i=1}^{I} \sum_{j \in A_i} (\Btot_{ij} - C_{ij}^{\mathrm{disc}}) \cdot x_{ij} & \text{(objective)} \\[6pt]
\text{Subject to:}\quad
& \sum_{i=1}^{I} \sum_{j \in A_i} C_{ij} \cdot x_{ij} \leq B^{\mathrm{total}} & \text{(total budget)} \\
& \sum_{j \in A_i} x_{ij} \leq 1, \qquad \forall i & \text{(mutual exclusivity)} \\
& x_{ij} \in \{0, 1\}, \qquad \forall i, \forall j \in A_i & \text{(binary)} \\[4pt]
& \sum_{i \in \mathcal{F}_m} \sum_{j \in A_i} C_{ij} \cdot x_{ij} \leq B_m, \quad \forall m & \text{(facility-type sub-budget)} \\
& \sum_{i \in \mathcal{R}_r} \sum_{j \in A_i} C_{ij} \cdot x_{ij} \leq B_r, \quad \forall r & \text{(regional cap)} \\
& \sum_{i \in \mathcal{R}_r} \sum_{j \in A_i} C_{ij} \cdot x_{ij} \geq \beta_r B^{\mathrm{total}}, \quad \forall r & \text{(regional floor)} \\[4pt]
& x_{ij} \leq x_{i'j'} & \text{(dependency, optional)} \\
& x_{ij} + x_{i'j'} \leq 1 & \text{(cross-project exclusivity, optional)} \\
& \sum_{j \in A_i} \Btot_{ij} \cdot x_{ij} \geq \theta \sum_{j \in A_i} C_{ij} \cdot x_{ij}, \quad \forall i & \text{(min BCR, optional)} \\
\end{aligned}
\end{equation*}

\noindent The cost decomposition $C_{ij} = \sum_{k=1}^{K_i} \sum_{n \in T_{jk}}
\mathit{Cost}_{ikn}$ (Equation~\ref{eq:cost-decomp}) is implicit in the
above; total benefit $\Btot_{ij}$ is computed per Equation~\ref{eq:btot}
with components from Sections~\ref{sec:safety}--\ref{sec:ccm}.

\paragraph{Problem scale.} The MILP has $\sum_{i=1}^{I} |A_i|$ binary
decision variables. The core-constraint count comprises one total-budget
constraint, $I$ mutual-exclusivity constraints, $|\mathcal{F}|$
facility-type sub-budgets, and up to $2|\mathcal{R}|$ regional constraints
(cap and floor). The optional constraints add at most one inequality per
declared dependency pair, conflict pair, or BCR-enforced project. For the
Harwood (2003) ten-site validation case, the problem has 41 binary
variables and 11 constraints; for the Banihashemi (2007) full network
case (135 homogeneous segments plus 13 intersections), the published
problem has 3{,}385 binary variables and 535 constraints.

% ---------------------------------------------------------------------
\subsection{Software Implementation}
% ---------------------------------------------------------------------

The MCBOMs MILP is implemented as a Python package, \texttt{mcboms}, that
constructs the model from input data and solves it via Gurobi
(commercial) or CBC (open-source) backends. The MILP itself --- the
mathematical statement summarized in
Section~\ref{sec:milp-summary} --- is solver-agnostic; it can be exported
to standard MILP exchange formats (AMPL \texttt{.mod}/\texttt{.dat},
GAMS \texttt{.gms}, CPLEX \texttt{.lp}, MPS) and solved by any compatible
commercial or open-source MILP solver, including IBM~CPLEX, FICO~Xpress,
Gurobi, SCIP, and CBC. Reference instances of the Harwood and
Banihashemi validation case studies are provided in these formats so
that any reviewer can independently reproduce the optimization results
without depending on the Python implementation.

The implementation is organized into a modular Python package,
\texttt{mcboms}, with the following structure (status as of Task~4
prototype):
\begin{itemize}[leftmargin=2em, itemsep=2pt]
  \item \texttt{mcboms.core} --- optimizer engine, alternative-set
    enumeration, and solver back-end interface (\emph{implemented}).
  \item \texttt{mcboms.benefits.safety} --- implementation of
    Equation~\ref{eq:safety-benefit-pwf} (\emph{implemented}, validated
    against the worked example to within floating-point precision and
    cross-validated against the parallel BCA spreadsheet to \$0.00 error
    across 35 segment--alternative pairs).
  \item \texttt{mcboms.benefits.operations} --- implementation of
    Equation~\ref{eq:ops-benefit} (\emph{in development}; the
    formulation is locked, the implementation will follow the same
    preprocessing pattern as \texttt{safety}).
  \item \texttt{mcboms.benefits.ccm} --- implementation of
    Equation~\ref{eq:ccm-benefit} (\emph{in development}; computational
    structure is locked per Section~\ref{sec:ccm}, category-specific
    monetization functions to be added incrementally).
  \item \texttt{mcboms.io} --- data I/O modules for HSM-format input,
    spreadsheet input, and validation case studies (\emph{implemented}
    for the BCA spreadsheet reader; HSM-format ingestion to follow).
  \item \texttt{mcboms.utils.economics} --- discount rates, present-worth
    factors, and unit crash cost defaults (USDOT BCA Guidance May 2025;
    \emph{implemented}).
  \item \texttt{mcboms.data} --- bundled validation case study data for
    Harwood (2003) and Banihashemi (2007) intersection sub-problem
    (\emph{implemented}).
\end{itemize}

The full source code is available in the project's GitHub repository,
along with an automated test suite (32 tests covering the worked example,
Harwood reproduction, and per-component unit tests), validation scripts
(\texttt{run\_harwood\_validation.py},
\texttt{banihashemi\_intersections.py}), and a programmer's guide.

% ---------------------------------------------------------------------
\subsection{Comparative Summary of Solution Methodologies}
% ---------------------------------------------------------------------

Task~3 evaluated four candidate solution methodologies for the MCBOMs
problem: MILP, surrogate-assisted genetic algorithms (GA), reinforcement
learning, and machine-learning regression on pre-computed benefit surfaces.
MILP was selected as the production methodology on the following bases:

\begin{itemize}[leftmargin=2em, itemsep=2pt]
  \item \textbf{Global optimality.} MILP solvers provide certified globally
    optimal solutions (within numerical tolerance), which is required for
    FHWA-deliverable resource allocation results.
  \item \textbf{Methodological transparency.} The objective function and
    constraints are explicit and auditable. Each step in the optimization
    can be traced back to the formulation, in contrast to black-box machine
    learning approaches.
  \item \textbf{Compatibility with HSM and CMF methodology.} The MILP
    formulation operates on benefit and cost values computed from
    HSM-standard methods (SPFs, CMFs, EB adjustment when applied),
    preserving the methodological lineage that FHWA deliverables require.
  \item \textbf{Maturity.} MILP for transportation resource allocation has
    a 30-year track record (Cohon, 1978; Harwood, 2003; Banihashemi, 2007;
    AASHTO Safety Analyst). Solver technology is mature.
  \item \textbf{Tractability.} The Harwood (2003) and Banihashemi (2007)
    case studies solve in under one second on commodity hardware. Realistic
    state-DOT problems with $10^4$--$10^5$ variables solve in minutes.
\end{itemize}

Machine learning methods are explicitly retained as preprocessing tools
where helpful (for example, gap-filling missing inputs to HSM SPFs), but are
not used as substitutes for the optimization core.

% ---------------------------------------------------------------------
\section{Validation Against Published Benchmarks}
\label{sec:validation}
% ---------------------------------------------------------------------

The MCBOMs implementation has been validated at three independent levels: a
methodology-level reproduction of the worked example
(Section~\ref{sec:safety-worked-example}), a program-level reproduction of
Harwood et al.~(2003), and a network-level reproduction of Banihashemi
(2007). Each validation level exercises a distinct portion of the framework
and tests against external published benchmarks rather than internal
self-consistency.

% ---------------------------------------------------------------------
\subsection{Methodology-Level Validation: Worked Example}
\label{sec:validation-worked}
% ---------------------------------------------------------------------

The worked example in Section~\ref{sec:safety-worked-example} computes a
single-segment safety benefit by direct application of
Equation~\ref{eq:safety-benefit-pwf}. The MCBOMs Python implementation
(\texttt{mcboms.benefits.safety.safety\_benefit}) reproduces this calculation
exactly: the annual safety benefit is computed as \$211{,}809, and the
20-year present-value benefit is \$2{,}243{,}913. The implementation
matches the manual calculation to within numerical floating-point precision
(maximum discrepancy less than \$0.01 across all severity components).

This validation confirms that the safety benefit module correctly implements
the methodology described in this document and that the framework default
unit crash costs are correctly applied.

% ---------------------------------------------------------------------
\subsection{Program-Level Validation: Harwood et al. (2003)}
\label{sec:validation-harwood}
% ---------------------------------------------------------------------

The Harwood et al.~(2003) study presents a 10-site case study covering a mix
of rural and urban undivided and divided nonfreeway facilities (2- and
4-lane), with optimization solved at two budget levels (\$10M and \$50M).
The MCBOMs MILP optimizer was run against the published Harwood case study
under the following configuration:

\begin{itemize}[leftmargin=2em, itemsep=2pt]
  \item Harwood's published per-site cost and benefit values from his
    Tables~2 and~3 (rather than recomputed from raw inputs)
  \item Harwood's discount rate of 4~percent (per AASHTO~1977)
  \item Harwood's two-bin severity convention (fatal+injury combined; PDO)
  \item No deferral-penalty mechanism (PNR is not implemented in the MCBOMs
    prototype; see the discussion in Section~\ref{sec:validation-harwood-10m}
    below)
\end{itemize}

\subsubsection{\$50 Million Budget Case}
\label{sec:validation-harwood-50m}

At the \$50M budget level, the budget constraint does not bind: all
cost-effective improvements can be funded simultaneously, regardless of any
deferral-penalty mechanism. Table~\ref{tab:harwood-50m} compares the MCBOMs
optimizer's results against the published Harwood values.

\begin{table}[H]
\centering
\caption{MCBOMs reproduction of Harwood (2003) Table~4, \$50~million budget.
Differences are attributable to rounding in the published values.}
\label{tab:harwood-50m}
\begin{tabular}{lrrr}
\toprule
\textbf{Metric} & \textbf{Harwood (2003)} & \textbf{MCBOMs} & \textbf{Difference} \\
\midrule
Total benefit (PSB + PTOB)             & \$10{,}640{,}914 & \$10{,}640{,}909 & $-\$5$ \\
Total expenditure (resurf + safety)    & \$16{,}271{,}246 & \$16{,}271{,}247 & $+\$1$ \\
Safety improvement cost                & \$4{,}481{,}397  & \$4{,}481{,}397  & \$0    \\
Net benefit                            & \$6{,}159{,}517  & \$6{,}159{,}512  & $-\$5$ \\
\midrule
Site selection (alternatives chosen)   & 10 sites match   & 10 sites match   & ---    \\
\bottomrule
\end{tabular}
\end{table}

The MCBOMs optimizer reproduces Harwood's published \$50M result to within
\$5 across all reported metrics, with identical site selection across all 10
sites. The five-dollar discrepancy is consistent with rounding in the
intermediate published values. This validates both the project-level MILP
formulation (Equations~\ref{eq:obj-project}--\ref{eq:binary-project}) and the
safety-benefit calculation (Equation~\ref{eq:safety-benefit-pwf}) at program
scale.

\subsubsection{\$10 Million Budget Case and the Treatment of Deferral
Penalties}
\label{sec:validation-harwood-10m}

At the \$10M budget level, the budget constraint binds: not all sites can be
fully funded, and some must be deferred (do-nothing alternative selected).
Harwood's published \$10M solution defers Sites~4, 6, and~9. The MCBOMs
optimizer, configured as described above, instead defers Sites~1, 4, 6,
and~8 and produces a higher net benefit. Table~\ref{tab:harwood-10m}
documents the comparison.

\begin{table}[H]
\centering
\caption{MCBOMs solution at \$10~million budget compared to Harwood (2003)
Table~4. The divergence reflects the deferral-penalty methodology described
below.}
\label{tab:harwood-10m}
\begin{tabular}{lrrr}
\toprule
\textbf{Metric} & \textbf{Harwood (2003)} & \textbf{MCBOMs} & \textbf{Difference} \\
\midrule
Total benefit (PSB + PTOB)             & \$7{,}187{,}814 & \$7{,}600{,}301 & $+5.7\%$ \\
Total expenditure (resurf + safety)    & \$9{,}953{,}579 & \$9{,}547{,}090 & $-4.1\%$ \\
Net benefit                            & \$4{,}675{,}033 & \$4{,}931{,}520 & $+5.5\%$ \\
\midrule
Sites deferred (do-nothing)            & \{4, 6, 9\}     & \{1, 4, 6, 8\}  & differs \\
\bottomrule
\end{tabular}
\end{table}

This divergence is methodologically intentional and reflects a documented
design choice. Harwood's full formulation includes a Penalty for Not
Resurfacing (PNR), which adds a deferral cost to the do-nothing alternative
roughly proportional to the avoided cost of eventual full pavement
reconstruction. The MCBOMs framework does \emph{not} implement PNR for the
following reasons:

\begin{itemize}[leftmargin=2em, itemsep=2pt]
  \item Harwood (2003) publishes only the aggregate PNR (\$5{,}576{,}145
    across the deferred sites in the \$10M case) and a qualitative formula.
    Per-site PNR values are not available, so reproducing Harwood's
    \$10M solution exactly would require reverse-engineering site-level
    penalties rather than computing them from a documented procedure.
  \item The PNR formula requires per-site pavement condition data and
    time-to-failure estimates that are not standardized in current FHWA
    guidance and are not present in the MCBOMs prototype's case study
    inputs.
  \item Banihashemi (2007), the second benchmark, also does not use a
    deferral-penalty mechanism.
\end{itemize}

Without PNR, the MCBOMs optimizer is free to defer sites strictly on
net-benefit grounds, and consequently finds a solution that is 5.5~percent
higher in net benefit than Harwood's published \$10M solution.
Both solutions are correct within their respective formulations: MCBOMs is
net-benefit-optimal subject to the budget constraint, while Harwood's
solution is optimal subject to budget plus PNR-incentivized resurfacing of
additional sites. Highway agencies wishing to incorporate pavement-deferral
costs in the optimization can do so by adding those costs to the
do-nothing alternative's cost term $C_{i,0}$ as a negative benefit; the
MCBOMs formulation supports this directly through
Equation~\ref{eq:obj-project}.

\subsubsection{Cross-Validation Against the Parallel BCA Spreadsheet}

Independent of the published-paper validations, the MCBOMs safety benefit
output has been cross-validated against the Task~4 data-curation BCA
spreadsheet on a per-(segment, alternative) basis. The spreadsheet
implements a streamlined formulation:
\begin{equation}
B^{\mathrm{safety,\ spreadsheet}}\ =\ \UC_K \cdot (\mathrm{Obs}_K - \mathrm{Proj}_K)
  + \UC_B \cdot (\mathrm{Obs}_{ABC} - \mathrm{Proj}_{ABC})
  + \UC_O \cdot (\mathrm{Obs}_O - \mathrm{Proj}_O)
\label{eq:spreadsheet-formula}
\end{equation}

\noindent which is Equation~\ref{eq:safety-benefit} reduced to (a)~a
single-year horizon (no discounting), and (b)~a three-bin severity
classification ($K$ / $ABC$-combined / $O$) with the Injury-B unit cost
applied to the entire $ABC$ bin. The cross-validation results across 35
segment--alternative pairs in the test dataset are:

\begin{itemize}[leftmargin=2em, itemsep=2pt]
  \item \textbf{Replication accuracy.} Maximum absolute error between the
    MCBOMs Python replication of Equation~\ref{eq:spreadsheet-formula} and
    the spreadsheet's reported \emph{Safety Benefit} values is \$0.00. The
    MCBOMs codebase reproduces every value in the spreadsheet's safety
    benefit column to the cent.
  \item \textbf{Methodological difference.} The ratio of MCBOMs
    Equation~\ref{eq:safety-benefit-pwf} present-value benefit to the
    spreadsheet's single-year benefit is constant at 10.59 across all
    35~pairs, exactly matching $\PWF(0.07, 20) = 10.594$. The two tools
    are mathematically equivalent up to the discounting choice and the
    severity-binning convention bridge.
\end{itemize}

% ---------------------------------------------------------------------
\subsection{Network-Level Validation: Banihashemi (2007)}
\label{sec:validation-banihashemi}
% ---------------------------------------------------------------------

The Banihashemi~(2007) study presents a network case study with three
highways (4--10 km in length) and 13 intersections, with the optimization
solved at ten budget levels from \$1M to \$12.2M. The case study uses the
IHSDM Crash Prediction Module (CPM) for both highway segment and
intersection crash predictions, with delay times computed from the Highway
Capacity Manual (Equations~A30-11 and~A30-15) for intersections and from
the published Highway~1 climbing-lane analysis for segments.

\paragraph{Validation scope.} The Banihashemi validation is conducted on the
intersection sub-problem of his case study: the 13 intersections with
published ADTs, improvement combinations, and delay times. The full network
reproduction (combining all three highways and all 13 intersections, totaling
135 homogeneous segments and 3{,}385 binary variables) requires
case-study-specific segment geometry data not published in the paper; that
full reproduction is documented as a planned extension.

\paragraph{Configuration.} The MCBOMs MILP optimizer was run against the
Banihashemi intersection sub-problem under the following configuration:
\begin{itemize}[leftmargin=2em, itemsep=2pt]
  \item IHSDM Crash Prediction Module (Equation~15 of Banihashemi),
    implemented with Banihashemi's published parameters ($\alpha$, $\beta$,
    $\gamma$, $C_i$) for 3-leg stop, 4-leg stop, and 4-leg signalized
    intersection types (page~113)
  \item Intersection type inferred from Table~3 traffic control column: all
    13 intersections treated as 4-leg stop-controlled or 4-leg signalized
    (Int~10) consistent with the traffic control entries
  \item AMF values for intersection features (skew, traffic control change,
    left-turn lane, sight distance deficiency) from IHSDM and Vogt/Bared
    base models; note that Banihashemi does not publish his exact AMF values,
    so exact numerical reproduction is not expected
  \item Unit delay cost $\$12.50$ per vehicle-hour (Banihashemi page~109)
  \item Average crash costs by facility type: $\$55{,}239$ (3-leg stop),
    $\$81{,}375$ (4-leg stop), $\$31{,}665$ (4-leg signalized)
    (Banihashemi page~109)
  \item 20-year project lifetime, no discounting (per Banihashemi's
    framework)
\end{itemize}

\paragraph{Results.} Table~\ref{tab:banihashemi-results} shows the MCBOMs
intersection selections at each budget level studied by Banihashemi.

\begin{table}[H]
\centering
\caption{MCBOMs intersection sub-problem results compared to the
intersection portion of Banihashemi (2007) Table~5. Notation: ``Int~$x$:
LTL'' means alternative with left-turn lane added at intersection $x$;
``SA'' means skew-angle correction; ``ISD'' means sight distance
improvement.}
\label{tab:banihashemi-results}
\scriptsize
\begin{tabular}{rp{5.4cm}p{5.4cm}}
\toprule
\textbf{Budget} & \textbf{MCBOMs (intersection sub-problem)} & \textbf{Banihashemi (full network)} \\
\midrule
Unlimited & Int~1:LTL; Int~2:LTL; Int~5:LTL; Int~6:LTL; Int~7:SA,ISD; Int~9:LTL; Int~12:LTL & Int~1:LTL; Int~2:LTL; Int~5:LTL; Int~6:LTL; Int~7:SA,ISD; Int~9:LTL; Int~10:LTL; Int~12:SA,LTL; Int~13:SA,LTL,ISD \\
\addlinespace
\$8M & same as unlimited (spent \$3.2M) & Int~1:LTL; Int~2:LTL; Int~5:LTL; Int~6:LTL; Int~12:LTL; Int~13:LTL \\
\addlinespace
\$6M & same as unlimited (spent \$3.2M) & Int~7:SA,ISD; Int~12:LTL; Int~13:LTL \\
\addlinespace
\$3M & Int~1:LTL; Int~2:LTL; Int~5:LTL; Int~6:LTL; Int~9:LTL; Int~12:LTL (spent \$2.1M) & Int~12:LTL; Int~13:LTL \\
\addlinespace
\$1M & Int~1:LTL; Int~2:LTL; Int~12:LTL (spent \$0.7M) & Int~12:LTL (spent within \$1M for intersections only) \\
\bottomrule
\end{tabular}
\end{table}

\paragraph{Interpretation of the comparison.} The MCBOMs optimizer and
Banihashemi's published solutions agree on three structural findings:

\begin{enumerate}[leftmargin=2em, itemsep=4pt]
  \item \textbf{The most cost-effective improvement is Int~12:LTL.} Both the
    MCBOMs optimizer and Banihashemi identify the left-turn lane at
    Intersection~12 as the highest-benefit-per-dollar improvement
    (benefit-to-cost ratio of approximately 11.5 in the MCBOMs computation;
    consistently selected at every non-zero budget in Banihashemi Table~5).
  \item \textbf{Signalization of Intersections 3 and 4 is never selected.}
    Both solvers correctly identify that signalizing these
    moderate-traffic minor-stop intersections would increase delay
    substantially (from 6{,}078 to 98{,}165 and 72{,}450 vehicle-hours/year
    respectively), making the delay cost increase dominate the crash cost
    reduction. Neither Banihashemi's solution nor MCBOMs selects these
    alternatives at any budget level.
  \item \textbf{Rank ordering of LTL improvements matches.} In both
    solutions, Intersections~12, 1, 2, 5, 6 appear as the first
    cost-effective LTL candidates as budget expands.
\end{enumerate}

The two solutions \emph{differ} in two respects that are attributable to
the scope of the sub-problem and to information not published in the
original paper:

\begin{enumerate}[leftmargin=2em, itemsep=4pt]
  \item \textbf{Budget coupling with segment improvements.} Banihashemi's
    Table~5 reports results for the full network; his intersection
    selections reflect the fact that segment improvements (lane widening,
    shoulder widening, etc.) are competing for the same budget. When
    isolated from the segment problem, the MCBOMs optimizer has the full
    budget available for intersections and selects more of them at any given
    budget level. This is a structural consequence of validating against a
    sub-problem rather than the full network.
  \item \textbf{AMF values not published.} Banihashemi does not publish the
    specific AMF values used for LTL, traffic control changes, or sight
    distance deficiency. The MCBOMs reproduction uses standard IHSDM and
    Vogt/Bared AMF values, which may differ from Banihashemi's choices by
    5--15~percent. This difference can shift the optimal alternative
    selection at the margin (for example, Intersection~10 LTL and several
    higher-cost alternatives at Intersections~12 and 13 that Banihashemi
    selects are marginal in the MCBOMs reproduction). Exact numerical
    reproduction of Banihashemi's selections at all budget levels would
    require his specific AMF values, which are not available.
\end{enumerate}

\paragraph{Validation conclusion.} The MCBOMs MILP framework correctly
implements the optimization structure used by Banihashemi~(2007), as
evidenced by (a)~consistent identification of the most cost-effective
improvement (Int~12:LTL) as the first-selected alternative in both
solutions, (b)~consistent rejection of signalization at Intersections~3
and~4 on delay-cost grounds, and (c)~consistent rank ordering of LTL
improvements across the intersection set. Residual divergence between the
two solutions is attributable to the intersection-sub-problem scope and to
Banihashemi's unpublished AMF values, neither of which constitutes a defect
in the MCBOMs formulation or implementation.

% ---------------------------------------------------------------------
\subsection{Summary of Validation Evidence}
% ---------------------------------------------------------------------

The MCBOMs framework has been validated at three independent levels:

\begin{enumerate}[leftmargin=2em, itemsep=4pt]
  \item \textbf{Methodology level} (Section~\ref{sec:validation-worked}):
    direct algebraic reproduction of the worked safety-benefit example.
    Match within floating-point precision.
  \item \textbf{Program level} (Section~\ref{sec:validation-harwood}): full
    optimization reproduction of the Harwood~(2003) \$50M published case to
    within \$5 (rounding) on all metrics, with identical site selection
    across all 10~sites. At the \$10M budget, the MCBOMs optimizer produces
    a strictly net-benefit-superior solution because the framework does not
    implement Harwood's PNR deferral-penalty mechanism; this methodological
    choice is documented and defensible.
  \item \textbf{Network level} (Section~\ref{sec:validation-banihashemi}):
    reproduction of the Banihashemi~(2007) intersection sub-problem. The
    MCBOMs MILP correctly identifies the most cost-effective intersection
    improvement (Int~12:LTL) as the first-selected alternative, correctly
    rejects signalization at Intersections~3 and~4 on delay-cost grounds,
    and produces a rank ordering of LTL improvements consistent with
    Banihashemi's published Table~5. Residual numerical divergence is
    attributable to the intersection-sub-problem scope (versus Banihashemi's
    full-network solution) and to AMF values not published in the original
    paper.
\end{enumerate}

\paragraph{Cross-validation.} In addition to the published-paper benchmarks,
per-segment cross-validation against the Task~4 data-curation BCA
spreadsheet shows perfect replication (\$0.00 maximum error across 35
segment--alternative pairs) of the spreadsheet's safety benefit calculation,
with the constant 10.59$\times$ ratio of present-value to single-year
benefit confirming methodological equivalence up to discounting.

\paragraph{Reproducibility.} All validation procedures described in this
document are reproducible from the project's GitHub repository. The
automated test suite includes 32~tests covering the worked example, the
Harwood reproduction, and unit tests for each safety-benefit computational
component. The Harwood validation script
(\texttt{run\_harwood\_validation.py}) prints a complete validation summary
on each run, including informational reporting of the documented \$10M
divergence.

% =====================================================================
% END MCBOMS METHODOLOGY CONTENT
% =====================================================================

\end{document}
